\scanpage{280}


Le \deleted{\unsure{cas}} cas le plus simple est celui où

\[
\Lambda=\mathbb Z/4\mathbb Z
\]

c'est aussi celui où $(\unsure{S_M|A_M},F_M)$ est justement l'octaèdre, dont on retrouve l'antipodisme par \deleted{voie} d'algèbre linéaire. La construction \inserted{\unsure{précédente|particulier}} de l'antipodisme est évidemment indépendante de l'hyp. \unsure{semi-locle|semi-locale} sur $\Lambda$, elle marche dès qu'on a $2\Lambda=\{0,2\}\ne\{0\}$, et dans $S_M$ on aura encore

\[
\left\{\begin{aligned}
&\operatorname{Et}_{\deleted{\unsure{u}}i}(u)=\operatorname{Et}_i(\underline a(u))\qquad\forall\,u\in S_M,\ i\in\unsure{\Pi_M|\pi_M}\\
&\underline a(u)\ne u
\end{aligned}\right.
\]

-- mais il n'est pas vrai nécessairement qu'il n'y a pas d'autres éléments que $u,u'\inserted{=\underline a(u)}$ ayant $\unsure{=|\text{m\textsuperscript{me}}}$ étoile que $u$ \ldots\ \deleted{Le cas dé}

\uline{NB} Dans tous les cas, on a

\[
\underline a\in\operatorname{Gl}'_\Lambda(M)^{\pm}
\]

car comme

\[
\underline a\equiv\mathrm{id}\quad\bmod 2\Lambda=\{0,2\unsure{]|\}},\quad\text{on a}
\]

\[
\det\underline a\in1+\{0,2\}=\{\pm1\}.
\]

En fait, on vérifie que

\[
\det\underline a=-1
\]

(car les éléments de $\unsure{N|\mathcal N}$ qui sont dans $\operatorname{Sl}_\Lambda(M)$ sont ceux de $\operatorname{Im}(M_0\oplus\mathbb F_2\,\mathrm{id})$, et $\underline a\in N\simeq\operatorname{End}_{\mathbb F_2}(M_0)$ n'est pas dans ce sous-module \ldots
% archive sheet number: 279
